%% =====================================================================
%%  CAPITULO 10 - CONCLUSAO
%% =====================================================================

\chapter{Conclusão}
\label{cap:conclusao}

\section{Resposta ao problema}
\label{sec:conc-resposta}

O problema investigado foi como estruturar um processo de desenvolvimento
assistido por agentes de inteligência artificial de modo que cada entrega fosse
rastreável a um requisito, verificável por um critério de aceitação e sustentada
por evidência reproduzível.

O trabalho responde em três partes.

\textbf{Quanto à verificabilidade, a resposta é afirmativa.} Um processo que
obriga a redigir critérios de aceitação antes da implementação, em forma que um
observador externo possa decidir o resultado, produz entregas verificáveis por
terceiros sem acesso ao raciocínio de quem implementou. Os doze critérios
declarados foram satisfeitos e verificados por observação externa ao produto.

\textbf{Quanto à sustentação por evidência, a resposta é afirmativa.} A
separação entre o que o agente afirma e o que a execução demonstra é o mecanismo
central do framework. Toda conclusão deste trabalho se apoia em inspeção textual
reproduzível, em execução observada ou em medição; nenhuma se apoia em afirmação
do executor.

\textbf{Quanto à rastreabilidade, a resposta é parcial.} A cadeia
requisito--critério--evidência existe e é verificável, mas cobre 54,2\,\% do
conjunto de requisitos. A cobertura é de 73,3\,\% para requisitos funcionais e de
apenas 22,2\,\% para não funcionais. O processo entrega rastreabilidade onde ela é
exigida com rigor, mas não garante cobertura completa, e a lacuna se concentra
exatamente na classe de requisito cuja verificação é mais difícil.

\section{Cumprimento dos objetivos}
\label{sec:conc-objetivos}

O Quadro~\ref{qua:cumprimento} retoma os objetivos específicos e registra o grau
de cumprimento.

\begin{quadro}[htb]
	\caption{Grau de cumprimento dos objetivos específicos}
	\label{qua:cumprimento}
	\centering
	\footnotesize
	\begin{tabularx}{\textwidth}{@{}p{1.1cm} X p{1.9cm}@{}}
		\toprule
		\textbf{ID} & \textbf{Resultado} & \textbf{Grau} \\
		\midrule
		OE-1 & Framework definido com cinco camadas, quinze agentes, quatro conjuntos de instruções, cinco habilidades, seis portões, catorze itens de conclusão e seis condições de parada. & Pleno \\
		\addlinespace
		OE-2 & Especificação com 15 requisitos funcionais, 9 não funcionais, 4 restrições, 12 critérios de aceitação e 5 casos de borda. & Pleno \\
		\addlinespace
		OE-3 & Sete decisões arquiteturais registradas com justificativa e alternativa descartada. & Pleno \\
		\addlinespace
		OE-4 & Artefato único de 1\,063 linhas, sem dependências externas, conforme as quatro restrições do alvo. & Pleno \\
		\addlinespace
		OE-5 & Doze de doze critérios satisfeitos, com evidência registrada por critério. & Pleno \\
		\addlinespace
		OE-6 & Sete ameaças à validade classificadas, seis dados ausentes marcados, lacuna de rastreabilidade identificada e localizada. & Parcial \\
		\bottomrule
	\end{tabularx}
	\legend{Fonte: elaborado pelo autor (2026).}
\end{quadro}

O objetivo específico OE-6 é classificado como parcial porque a avaliação
identificou ameaças que não puderam ser mitigadas dentro do escopo, em particular
a ausência de grupo de controle e a ausência de repetição da medição. Identificar
e declarar uma limitação é resultado válido, mas não é o mesmo que eliminá-la.

\section{Contribuições}
\label{sec:conc-contribuicoes}

\begin{enumerate}
	\item \textbf{Um protocolo de controle operacional} para desenvolvimento com
	agentes de inteligência artificial, com precedência do critério sobre a
	implementação e proibição de evidência fabricada, materializado em artefatos
	versionados e reutilizáveis.

	\item \textbf{Um método de verificação por observação externa} que não
	instrumenta o produto sob teste, viável em interfaces gráficas discretas por
	amostragem da superfície de desenho e uso de agente de teste autônomo.

	\item \textbf{Evidência empírica de uma assimetria} entre a verificabilidade
	de requisitos funcionais e não funcionais, com medição da cobertura por classe:
	73,3\,\% contra 22,2\,\%.

	\item \textbf{Demonstração de que casos de borda declarados alteram o projeto},
	e não apenas os testes, com um caso concreto em que a declaração prévia criou um
	estado de máquina que não existia na modelagem inicial.

	\item \textbf{Registro do custo e das limitações} do processo, incluindo a
	declaração explícita do que não foi verificado e de quais dados não foram
	obtidos, em vez de preenchimento por estimativa.
\end{enumerate}

\section{Limitações}
\label{sec:conc-limitacoes}

As limitações deste trabalho são de quatro naturezas.

\textbf{Metodológica.} Não houve grupo de controle; o desenho é de estudo de caso
único, com generalização analítica e não estatística. Não houve repetição da
medição de intervalo, e portanto não há inferência estatística sobre ela.

\textbf{De escopo.} O alvo tem estado discreto, regras determinísticas e ausência
de concorrência, integração externa e dependência de hardware. Os resultados não
se transferem diretamente a domínios de tempo real, sistemas distribuídos ou
software embarcado, nos quais o framework de origem previa níveis de evidência
superiores.

\textbf{De cobertura.} Sete dos nove requisitos não funcionais não possuem
critério de aceitação. Sete itens permaneceram fora da verificação, incluindo
validação em dispositivo de toque físico, compatibilidade entre navegadores e
comportamento sob armazenamento local indisponível.

\textbf{De evidência.} Nenhum dado de consumo de memória, de tempo de execução do
processo ou de custo por tarefa foi coletado. As afirmações sobre custo do
processo são qualitativas.

\section{Trabalhos futuros}
\label{sec:conc-futuros}

\begin{enumerate}
	\item \textbf{Adicionar portão de requisito órfão.} Estender o portão G1 para
	rejeitar requisito sem critério de aceitação, ou obrigar seu rebaixamento
	explícito a diretriz não verificável. É a correção diretamente indicada pelo
	achado principal deste trabalho.

	\item \textbf{Instrumentar o processo.} Medir tempo, número de iterações e
	custo de contexto por tarefa, convertendo a avaliação qualitativa de custo em
	avaliação quantitativa.

	\item \textbf{Executar experimento controlado.} Desenvolver o mesmo alvo com e
	sem o processo, com métricas comparáveis, para atacar a ameaça AV-1.

	\item \textbf{Repetir a medição de velocidade} com múltiplas execuções e
	tratamento estatístico, com registro do número de amostras por faixa.

	\item \textbf{Aplicar a domínio de nível de evidência superior.} Verificar o
	comportamento do protocolo em desenvolvimento embarcado, com teste em bancada e
	integração com hardware, domínio para o qual o material de origem do framework
	foi concebido.

	\item \textbf{Avaliar a conversão de requisitos não funcionais em métricas}
	como etapa assistida por agente, verificando se um agente dedicado reduz a
	assimetria medida.
\end{enumerate}

\section{Considerações finais}
\label{sec:conc-finais}

A contribuição central deste trabalho não é um produto de software nem um modelo
de linguagem. É uma demonstração de que a disciplina de especificação, aplicada a
um fluxo conduzido por agentes de inteligência artificial, produz um resultado
auditável --- e de que essa disciplina tem um ponto cego mensurável exatamente
onde a verificação é mais difícil.

O achado de que a cobertura de requisitos não funcionais é de 22,2\,\% é, do
ponto de vista da engenharia, mais útil que a afirmação de que doze critérios
foram satisfeitos. O primeiro resultado localiza um defeito no processo e indica
a correção; o segundo apenas confirma que o processo funcionou onde era fácil
funcionar. Um processo de verificação que não expõe sua própria lacuna é um
processo que ainda não foi verificado.
